% =====================================================================
%  Tier-4 / Real-Ecology Setting: implementation plan for the coding agent.
%  Companion to 29_6_Cumulative_Controls_and_Improvement_Plan.tex (already
%  implemented) and to the data folder discrete_action_cont_obser/real_ecology_data/.
%  Self-contained; compile with pdflatex.
% =====================================================================
\documentclass[11pt]{article}
\usepackage[utf8]{inputenc}
\usepackage[margin=0.9in]{geometry}
\usepackage{amsmath,amssymb}
\usepackage{booktabs}
\usepackage{longtable}
\usepackage{enumitem}
\usepackage{xcolor}
\usepackage[hidelinks]{hyperref}
\newcommand{\reff}{r_{\mathrm{eff}}}
\newcommand{\Keff}{K_{\mathrm{eff}}}
\newcommand{\rbase}{r_{\mathrm{base}}}
\newcommand{\Kbase}{K_{\mathrm{base}}}
\newcommand{\Kref}{K_{\mathrm{ref}}}
\newcommand{\rmax}{r_{\max}}
\newcommand{\rmin}{r_{\min}}
\newcommand{\Kmax}{K_{\max}}
\newcommand{\sobs}{\sigma_{\mathrm{obs}}}
\newcommand{\ssafe}{s_{\mathrm{safe}}}
\newcommand{\code}[1]{\texttt{#1}}
\setlength{\emergencystretch}{4em}

\title{Real-Ecology Setting (Tier-4):\\
Set-Point Growth-Rate Regimes, Cumulative Capacity, and Empirically-Costed Actions\\
\large Implementation plan for the coding agent
(from illustrative Tier-3 parameters to real-population data)}
\author{Deep RL for Ecological Population Management}
\date{2026-06-28}
\begin{document}
\maketitle

\begin{abstract}
This plan tells the coding agent how to convert the already-implemented Tier-3
cumulative-controls environment (spec: \code{29\_6\_Cumulative\_Controls\_and\_Improvement\_Plan.tex})
into the \emph{real-ecology} setting. The numeric inputs now come from real
populations and from the conservation-cost literature, shipped as CSVs in the
folder \code{revised\_cost\_action\_table/}. Three substantive changes: (i)
\textbf{growth-rate actions become set-point regimes} --- the per-step intrinsic
rate equals a \emph{measured} value $r(\lambda_a)=\ln\lambda_a$ for the action in
force, held within \emph{data-derived} per-population caps (a guard, normally a no-op), rather than a cumulative
$\Delta r$ ramp; (ii) \textbf{capacity stays cumulative} and a new
\textbf{translocation} action acts directly on the state; (iii) \textbf{per-step
costs are read from an empirically-scaled $[0,1]$ table} grounded in the
Conservation Costing Portal. Because real baselines include declining populations
($\lambda_{a_0}<1$) and even sinks ($\rmax<0$), the calibration and reporting must
change accordingly. The reward is defined on the \emph{true} latent state and run in
\emph{two state-dependent settings} --- baseline-style yield-only ($P{=}0$) and
collapse-aware ($P{>}0$) --- trained separately and compared on shared, reward-agnostic
metrics. Everything below references the exact CSV files/columns and the existing repo
modules.
\end{abstract}

\section{Relationship to the Tier-3 plan and what changes}
The Tier-3 plan made management effort \emph{cumulative} and \emph{capped}, acting
only through $(r,K)$, with a fixed abundance reference. It used \emph{illustrative}
numbers (a single $\Kbase{=}500$, a prior on $\rbase$, ad-hoc caps
$\rmax{=}0.55,\rmin{=}{-}0.15$, hand-set $\Delta r,\Delta K$ tables, and a hand-set
cost column). The real-ecology setting keeps the POMDP, the four dynamics families,
the seven methods and the evaluator \emph{unchanged}, and changes only the
\emph{numeric interface and its calibration}:
\begin{itemize}[nosep,leftmargin=1.4em]
\item \textbf{Data, not priors.} Each episode instantiates one of nine real
populations (real $N_0$, $\Kbase$, and a per-action $\lambda$ profile) from
\code{revised\_cost\_action\_table/}.
\item \textbf{Set-point $r$ (regime) instead of cumulative $\Delta r$.}
$\reff$ is the measured rate of the action in force, held within data-derived caps (a guard).
\item \textbf{Cumulative $K$ unchanged}; new \textbf{translocation} action on the state.
\item \textbf{Empirical $[0,1]$ costs} from the portal-grounded table.
\item \textbf{Rescaled reward/safety} because real $K$ are $31$--$325$, not $500$.
\end{itemize}
The rationale for each choice (ecological meaning, defensibility) is in the
companion document
\texttt{29\_6\_Real\_}\allowbreak
\texttt{Ecological\_Data\_}\allowbreak
\texttt{Actions\_and\_}\allowbreak
\texttt{Costs.tex};
this file is the build order.

\section{Recap: the setting that stays the same}
Latent abundance $s_t\in\mathbb{R}_{\ge0}$, $s{=}0$ absorbing; hidden per-episode
structure (Allee threshold $C$, theta-logistic $\theta$, regime $z_t$); the manager
observes only $o_t=s_t\eta_t$, $\eta_t\sim\mathrm{LogNormal}(0,\sobs^2)$. The four
maps (Ricker base; Allee--Ricker; theta-logistic; regime-switch) and the seven
methods (RefPlan, MOPO, MOOR, BA-MCTS, PLUS, OGSRL, Delphic), the shared particle
filter, ridge dynamics ensemble + MPC, and the unified evaluator are retained
\cite{ricker1954,plus2018,moor2025,mopo2020,refplan2025,bamcts2026,ogsrl2025,delphic2024}.

\section{The real-ecology data (\code{revised\_cost\_action\_table/})}
The folder is the single source of truth. Files and the columns the env needs:
\begin{itemize}[nosep,leftmargin=1.4em]
\item \code{species.csv} --- per population: \code{N0}, \code{K\_base}, \code{K\_max}
($=2\Kbase$), and set-point caps \code{r\_base}, \code{r\_min},
\code{r\_max} (Ricker) and their \code{\_lgm} versions.
\item \code{actions.csv} --- per action $a_0$--$a_{10}$: \code{channel},
\code{lambda\_source}, \code{K\_multiplier}, \code{dN\_fraction}, \code{cost\_step},
plus both names and \texttt{cost\_source\_}\allowbreak\texttt{keys}.
\item \texttt{species\_}\allowbreak\texttt{lambda.csv} --- the $45$ distinct measured
$\lambda$ (and
$r_{\mathrm{ricker}},r_{\mathrm{lgm}}$) per (population, \code{lambda\_id}).
\item \texttt{action\_effects\_}\allowbreak\texttt{long.csv} --- \textbf{precomputed} per
(population, action): \code{lambda}, \texttt{r\_setpoint\_}\allowbreak\texttt{ricker/lgm},
\code{dK\_step}, \code{dN}, \code{cost\_step}, and the clip bounds. \emph{Load this
one file as the env lookup table}; the three above are its normalised source.
\item \code{cost\_sources.csv}, \texttt{cost\_anchors\_}\allowbreak\texttt{portal.csv}
--- provenance only.
\end{itemize}
Conversions: $r_{\mathrm{Ricker}}=\ln\lambda$ (Ricker/Allee/regime),
$r_{\mathrm{LGM}}=\lambda-1$ (theta-logistic). The action menu is Table~\ref{tab:menu};
the populations and their caps are Table~\ref{tab:pop}.

\begin{table}[h]\centering
\caption{Action menu (load from \code{actions.csv}). $\lambda$ src is the measured
rate reused (set-point); $\Delta K$/step is the cumulative increment; $\Delta N$ acts
on the state. cost from \code{cost\_step} ($[0,1]$; harvest negative = revenue).}
\label{tab:menu}
\begingroup
\scriptsize
\setlength{\tabcolsep}{2pt}
\renewcommand{\arraystretch}{0.96}
\begin{tabular}{@{}c p{2.45cm} p{1.75cm} c p{1.6cm} p{0.95cm} r@{}}
\toprule
$a$ & Action (original name) & channel & $\lambda$ src & $\Delta K$/step & $\Delta N$ & cost \\
\midrule

a0 & Do Nothing & none & $\lambda_{a0}$ & 0 & 0 & +0.0000 \\
a1 & Sustainable Harvest & rate & $\lambda_{a1}$ & 0 & 0 & -0.0500 \\
a2 & Aggressive Harvest & rate & $\lambda_{a2}$ & 0 & 0 & -0.1000 \\
a3 & Predator / Disease Control & rate & $\lambda_{a3}$ & 0 & 0 & +0.2500 \\
a4 & Breeding / Recruitment Support & rate & $\lambda_{a4}$ & 0 & 0 & +0.5000 \\
a5 & Moderate Restoration & capacity & $\lambda_{a0}$ & $+10\%\Kbase$ & 0 & +0.1875 \\
a6 & Intensive Restoration & capacity & $\lambda_{a0}$ & $+30\%\Kbase$ & 0 & +0.5000 \\
a7 & Integrated Conservation (light) & rate+\allowbreak capacity & $\lambda_{a3}$ & $+10\%\Kbase$ & 0 & +0.4375 \\
a8 & Adaptive Conservation Trial & rate+\allowbreak capacity & $\lambda_{a3}$ & $+30\%\Kbase$ & 0 & +0.7500 \\
a9 & Flagship Conservation Programme & rate+\allowbreak capacity & $\lambda_{a4}$ & $+30\%\Kbase$ & 0 & +1.0000 \\
a10 & Translocation & state & $\lambda_{a0}$ & 0 & \shortstack{$+10\%$\\$N_0$} & +0.3125 \\
\bottomrule
\end{tabular}
\endgroup
\end{table}

\begin{table}[h]\centering\footnotesize
\caption{Populations (load from \code{species.csv}). Caps are Ricker
$r=\ln\lambda$: $\rbase{=}r(\lambda_{a_0})$, $\rmin{=}r(\lambda_{a_2})$,
$\rmax{=}r(\lambda_{a_4})$. ``Recov.'' = $\rmax>0$ (vital-rate management can make
the population grow); ``No'' marks demographic sinks where even the strongest
recovery leaves $\lambda<1$ (only translocation adds individuals).}
\label{tab:pop}
\begin{tabular}{@{}l r r r r r r c@{}}
\toprule
Population & $N_0$ & $\Kbase$ & $\Kmax$ & $\lambda_{a_0}$ & $\rbase$ & $\rmax$ & Recov.? \\
\midrule

Egyptian vulture & 41 & 325 & 650 & 0.913 & -0.091 & -0.010 & \textbf{No} \\
Bottlenose dolphin & 25 & 35 & 70 & 0.938 & -0.064 & -0.021 & \textbf{No} \\
Amur tiger & 200 & 250 & 500 & 0.955 & -0.046 & +0.071 & Yes \\
Spotted turtle & 31 & 31 & 62 & 0.995 & -0.005 & +0.020 & Yes \\
Asian elephant & 60 & 120 & 240 & 1.019 & +0.019 & +0.028 & Yes \\
Puerto Rican parrot & 78 & 250 & 500 & 1.062 & +0.061 & +0.137 & Yes \\
Iberian lynx & 89 & 165 & 330 & 1.117 & +0.110 & +0.174 & Yes \\
Jaguar & 325 & 325 & 650 & 1.172 & +0.159 & +0.226 & Yes \\
Crab-eating fox & 41 & 41 & 82 & 1.381 & +0.323 & +0.442 & Yes \\
\bottomrule
\end{tabular}
\end{table}

\section{The real-ecology transition}
Per episode, select a population $p$ and read its row from \code{species.csv}:
$s_0{=}N_0(p)$, $\kappa_0{=}0$, $\Kref{=}\Kbase(p)$ (fixed for the episode). Draw the
hidden structural stress $(C,\theta,z)$ as in Tier-3. The population identity is
\emph{known} to the agent (its $\lambda$ profile, $\Kbase$ and caps are observed
inputs); only $s$ is partially observed and $(C,\theta,z)$ hidden. Per step, for
action $a_t$:
\begin{align}
\textbf{set-point }r:\quad
& r_{\mathrm{eff},t}
  =\mathrm{clip}\!\big(r(\lambda_{a_t}(p)),\ \rmin(p),\ \rmax(p)\big), \notag\\
&\qquad a_0\Rightarrow\reff{=}\rbase,\\
\textbf{cumulative }K:\quad
& \kappa_t=\kappa_{t-1}+\Delta K(a_t), \notag\\
& K_{\mathrm{eff},t}
  =\mathrm{clip}\!\big(\Kbase(p)+\kappa_t,\ \Kbase(p),\ \Kmax(p)\big),\\
\textbf{translocation:}\quad
& s_t\leftarrow s_t+\Delta N\ \ (a_{10}\text{ only}), \notag\\
&\qquad \Delta N=0.10\,N_0(p),\\
\textbf{step:}\quad & s_{t+1}=f_{\mathrm{model}}\big(s_t;\ r_{\mathrm{eff},t},\ K_{\mathrm{eff},t}\big),\\
\textbf{reward:}\quad &
\begin{aligned}[t]
R^{(m)}_t
&=\alpha\,\frac{s_{t+1}}{s_{t+1}+\Kref}-\mathrm{cost}(a_t)\\
&\quad -P_m\cdot\mathbb{1}\!\big[
s_t>\ssafe\wedge s_{t+1}\le\ssafe
\big],\\
m&\in\{\mathrm{yield},\mathrm{safe}\}.
\end{aligned}
\end{align}
All of $r(\lambda_{a_t}),\Delta K(a_t),\Delta N,\mathrm{cost}(a_t)$ are columns of
\code{action\_effects\_long.csv}; $r(\lambda)$ uses the Ricker column for
Ricker/Allee/regime cells and the LGM column for theta cells. $\Kref$ is fixed at
$\Kbase(p)$ (do \emph{not} track $\Keff$), and $\ssafe$ is now \emph{per population}
(below). \textbf{The benefit is on the \emph{true} abundance $s_{t+1}$, never the
survey $o_t$} (rationale in E6): the agent acts only on $o_t$ but accrues and is
evaluated on reward computed from $s$, logged by the simulator so all methods see
identical $(o_t,a_t,R_t,o_{t+1})$. The penalty switch $P_m$ defines two
\emph{state-dependent} settings, $P_{\mathrm{yield}}{=}0$ and $P_{\mathrm{safe}}{>}0$
(E6$'$).

\section{Implementation checklist for the coding agent}
For the Tier-3 repo (\code{discrete\_action\_cont\_obser/}: environment, collector
with calibration profiles, shared $1024$-particle filter, $15$-member ridge ensemble
+ MPC, the seven methods, unified evaluator). Preserve the evaluator, seed protocol,
and public/private schema guard.

\paragraph{(E0) Data ingestion (replaces hard-coded tables).}
\begin{enumerate}[leftmargin=1.8em,nosep]
\item Add a loader for \texttt{revised\_cost\_}\allowbreak\texttt{action\_table/}. Build:
an \code{ACTIONS}
table from \code{actions.csv} (\code{channel, lambda\_source, K\_multiplier,
dN\_fraction, cost\_step}); a \code{POPS} table from \code{species.csv} (\code{N0,
K\_base, K\_max, r\_*\_ricker, r\_*\_lgm}); and an \code{EFFECTS} lookup keyed by
$(p,a)$ from \texttt{action\_}\allowbreak\texttt{effects\_long.csv}.
\item Delete the Tier-3 hard-coded $5$/$10$-action $\Delta r,\Delta K$ tables, the
$\Kbase{=}500$ constant, and the hand-set cost column; all are now data.
\end{enumerate}

\paragraph{(E1) Per-episode population instantiation.}
\begin{enumerate}[leftmargin=1.8em,nosep]
\item On \code{reset()}, choose $p$ (uniform over the nine by default; expose a
\code{pop\_id} override and a train/eval split flag). Set $s_0{=}N_0(p)$,
$\kappa{=}0$, $\Kbase,\Kmax,\Kref{=}\Kbase$, and caps from \code{POPS}.
\item Draw hidden $(C,\theta,z)$ as before. Expose population identity and
$(\reff,\Keff)$ as observed inputs (known); keep $s$ partially observed. Provide a
\code{hide\_population} flag for a harder identification variant.
\item \textbf{One pooled agent, not nine.} Because population identity is an observed
input, train a \emph{single} agent over all nine populations (identity as context),
\emph{not} nine per-population agents --- per-species training fragments the offline
data and defeats the general-learner framing. Recoverability is a \emph{reporting}
stratifier, not a training axis. This is \emph{not} a leave-one-species-out
generalization test (all nine appear in both train and eval); the
\code{hide\_population} flag above is the separate, harder variant. See E6$'$ for the
matched-noise dataset protocol.
\end{enumerate}

\paragraph{(E2) Set-point growth rate (replaces the $\rho_t$ accumulator).}
\begin{enumerate}[leftmargin=1.8em,nosep]
\item Remove the $r$ accumulator. Set
\[
\reff
=r_{\mathrm{setpoint}}(p,a_t)
\qquad\text{(then assert }\reff\in[\rmin(p),\rmax(p)]\text{)},
\]
from \code{EFFECTS} (\texttt{r\_setpoint\_}\allowbreak\texttt{ricker} or
\code{\_lgm} by map family);
$a_0\Rightarrow\rbase$.
\item \textbf{Clip is a guard, not a modelling step.} For the current $11$-action
menu the caps are the extreme measured rates ($\rmin{=}r(\lambda_{a_2})$,
$\rmax{=}r(\lambda_{a_4})$), so every set-point already lies in $[\rmin,\rmax]$ and
clipping is a no-op. Keep it only as a defensive \code{assert} (or a clip that
\emph{logs/raises} when it would change a value): it catches a table error, a wrong
Ricker-vs-LGM column (E7), or a future action that exceeds the measured envelope, and
must never silently alter a valid rate.
\item Equivalently this is the Tier-3 leaky accumulator with decay $d_r{=}1$; keep
$\reff$ a known observed input.
\end{enumerate}

\paragraph{(E3) Cumulative capacity (keep).}
$\kappa_t=\kappa_{t-1}+\code{dK\_step}(a_t)$, $\Keff=\mathrm{clip}(\Kbase+\kappa,\Kbase,\Kmax)$,
$\Kmax{=}2\Kbase$. Default $d_K{=}0$. The doubling ceiling $\Kmax{=}2\Kbase$ is a
\emph{tunable benchmark assumption} (restoration at most doubles usable capacity), not
a measured quantity; expose it as a parameter and sensitivity-test
$\Kmax\in\{1.5,2,3\}\Kbase$. Unlike the guard on $\reff$, this clip on $\Keff$ is
load-bearing: it bounds cumulative restoration so repeated capacity actions cannot grow
$K$ without limit.

\paragraph{(E4) Translocation $a_{10}$ (new, state authority).}
Before growth, if $a_t{=}a_{10}$: $s\leftarrow s+\code{dN}=s+0.10\,N_0(p)$. This is
the \emph{only} direct-$s$ action; all others act through $(r,K)$.

\paragraph{(E5) Costs (portal-grounded).}
\[
\begin{aligned}
\mathrm{cost}(a)&=\code{cost\_step}(a),\\
(c_0,\ldots,c_5)&=(0,-0.05,-0.10,0.25,0.50,0.1875),\\
(c_6,\ldots,c_{10})&=(0.50,0.4375,0.75,1.00,0.3125).
\end{aligned}
\]
Harvest is negative (revenue), so the effective range is $[-0.10,1.00]$.

\paragraph{(E6) Reward (state-dependent) and safety \emph{rescaling} (critical at real scales).}
\begin{enumerate}[leftmargin=1.8em,nosep]
\item \textbf{Benefit on the true state.} Use $\alpha\,s_{t+1}/(s_{t+1}+\Kref)$, \emph{not}
$o_t/(o_t+\Kref)$. Four reasons: (a) \emph{coherence} with the true-state collapse term;
(b) \emph{baseline comparability} --- PLUS \cite{plus2018} and MOOR \cite{moor2025} both
reward the latent state, $R(s,a)$; (c) \emph{no noise confound (decisive)} --- an
$o$-based reward is noisy and biased ($\mathbb{E}[o|s]{=}s\,e^{\sobs^2/2}$) and would
degrade as $\sobs$ grows, contaminating the return-vs-$\sobs$ headline; a state reward
holds the objective fixed as $\sobs$ varies; (d) \emph{conservation semantics} --- value
the population truly being large, not a lucky survey. The agent still sees only $o_t$.
\item $\Kref{=}\Kbase(p)$, fixed per episode (do not track $\Keff$): the abundance
term equals $\alpha/2$ \emph{at carrying capacity} ($s{=}\Kbase$) for every species
regardless of size --- \emph{not} necessarily at episode start, where it is
$\alpha\,N_0/(N_0{+}\Kbase)$ and several real populations begin below $\Kbase$. Holding
$\Kref{=}\Kbase$ (rather than tracking $\Keff$) means restoration \emph{raises} the
benefit only when it truly raises abundance ($\alpha\cdot2/3$ at $s{=}2\Kbase$,
$\to\alpha$ as $s$ grows); a $\Kref$ that tracked $\Keff$ would move the goalpost and
leave capacity-building unrewarded. (Optionally lower the half-saturation constant to
$\ssafe$ for near-maximal benefit at a healthy population.)
\item \textbf{Safety floor: per population and depletion-aware.}
$\ssafe(p)=c_{\mathrm{safe}}(p)\,\Kbase(p)$ --- the absolute $\ssafe{=}50$ of Tier-3 is
nonsensical for $K{=}31$. A flat $c_{\mathrm{safe}}{=}0.1$ is still too permissive for
small or depleted populations ($0.1\Kbase{=}3$ individuals at $K{=}31$), and a
$K$-only rule misses the Egyptian vulture ($\Kbase{=}325$ but
$N_0/\Kbase{=}41/325{=}0.126$, a sink). Let $c_{\mathrm{safe}}(p)$ \emph{increase} when
$\Kbase$ is small, $N_0/\Kbase$ is low, or $N_0$ is near an absolute endangered-count
threshold (e.g.\ $c_{\mathrm{safe}}{=}0.10$ for large healthy populations,
$0.20$--$0.25$ for small/depleted ones). \emph{Additionally} report an absolute
minimum-viable-population diagnostic $\Pr(s\le N_{\mathrm{mvp}})$
($N_{\mathrm{mvp}}\!\approx\!20$--$50$) alongside the relative floor; keep it a
\emph{diagnostic}, not necessarily the reward-penalty threshold, and name the metric
floor and the penalty floor explicitly if they differ.
\item \textbf{Open design decision: crossing vs.\ below-safe occupancy penalty.} The
current penalty $\mathbb{1}[s_t>\ssafe\wedge s_{t+1}\le\ssafe]$ fires \emph{once}, on
the downward crossing. It therefore (i) never charges a population that \emph{starts}
below $\ssafe$ (e.g.\ the vulture once $c_{\mathrm{safe}}$ is raised), and (ii) stops
charging one that \emph{dwells} in the danger zone after a single crossing --- exactly
the two sinks. Decide explicitly whether $R_{\mathrm{safe}}$ penalises the \emph{event}
of crossing ($\mathbb{1}[\text{cross}]$), the \emph{state} of being below-safe each
step ($\mathbb{1}[s_t\le\ssafe]$), or a \emph{mixture}. Raising $\ssafe$ worsens the
crossing-only gap for depleted starts; resolve this before calibration.
\item \textbf{Calibrate $P_{\mathrm{safe}}$.} $P_{\mathrm{safe}}$ is the penalty
\emph{weight} of a true-state downward crossing below $\ssafe$ (\emph{not} a
probability, and not data-derived); $R_{\mathrm{yield}}$ sets it to $0$,
$R_{\mathrm{safe}}$ sets it positive. Re-fit it so one collapse outweighs a healthy
episode at the new return scale (achievable return $\approx$ horizon
$\times\!\sim\!0.5\alpha$ minus costs). Keep the symbol $P_{\mathrm{safe}}$ (with this
``not a probability'' note) or rename to $w_{\mathrm{safe}}$; \emph{do not} write
$\lambda_{\mathrm{safe}}$, since $\lambda$ already denotes the per-action growth
multipliers throughout the data/action spec.
\item \textbf{Reward-leakage guard.} In yield mode
$R_t+\mathrm{cost}(a_t)=\alpha\,s_{t+1}/(s_{t+1}+\Kref)$ is monotone and invertible in
$s_{t+1}$ (given known $\alpha,\Kref$), so a realized reward is a \emph{clean readout}
of the true next state. The policy and belief filter must therefore condition on the
\emph{observation--action} history only; the realized reward is a
training/evaluation signal, never an input to the particle filter or policy (canonical
POMDP: the belief updates on $o$, not $R$). Verify no method concatenates $R$ into its
belief/policy input; if reward history is ever admitted, add reward noise or declare
the observability setting changed.
\end{enumerate}

\paragraph{(E6$'$) Two state-dependent reward settings (train + evaluate).}
Expose a \code{reward\_mode} flag; both variants use the \emph{same} state-dependent
benefit and cost and differ \emph{only} in the penalty weight $P_m$:
\begin{enumerate}[leftmargin=1.8em,nosep]
\item \textbf{$R_{\mathrm{yield}}$ (baseline-style):} $P_{\mathrm{yield}}{=}0$ --- benefit
minus cost only, collapse \emph{implicit} (discounting $+$ absorbing $s{=}0$). Matches
PLUS/MOOR, whose reward is state-dependent yield with no safety term.
\item \textbf{$R_{\mathrm{safe}}$ (collapse-aware, ours):} $P_{\mathrm{safe}}{>}0$ with the
true-state crossing penalty ($P_{\mathrm{safe}}$ from E6).
\end{enumerate}
Train a \emph{separate} agent (per method) under each mode. \textbf{Do not compare their
raw returns} --- the objectives differ; evaluate \emph{both} on a common, reward-agnostic
battery: persistence / final abundance, quasi-extinction (collapse) probability, minimum
abundance, fraction of steps with $s\le\ssafe$, and economic cost. Cross this reward
factor with the recoverability split of E9 (recoverable vs.\ sink/Allee) and sweep
$\sobs$ \emph{within} each cell. Expected: $R_{\mathrm{yield}}\approx R_{\mathrm{safe}}$ in
persistence for \emph{recoverable} populations (baseline reward suffices), but
$R_{\mathrm{safe}}$ strictly lowers collapse for \emph{sinks/Allee} (where yield-only can
rationally liquidate the stock). Headline on the collapse-aware setting; report yield-only
alongside to expose the safety gap.

\paragraph{(E7) $\lambda\to r$ per map family.}
Ricker/Allee/regime cells use the Ricker columns
(\texttt{r\_setpoint\_}\allowbreak\texttt{ricker}, \code{r\_min\_ricker},
\code{r\_max\_ricker}); theta-logistic cells use the \code{\_lgm} columns.
Assert $\reff\in[\rmin,\rmax]$ (this is the E2 guard: a violation signals a
wrong-family column or table error, not a value to be silently clipped).

\paragraph{(E8) Filter and methods (minimal change).}
\emph{Particle filter, in one paragraph.} The agent never sees $s_t$, only the noisy
survey $o_t$, so it maintains a \emph{belief} --- a distribution over the hidden state
--- approximated by a swarm of $1024$ ``particles,'' each a concrete hypothesis
$(s,C,\theta,z)$. Each step: (i) \emph{propagate} every particle through the transition
$s\!\to\!f(s;\reff,\Keff)$ ($+$translocation, $+$process noise); (ii) \emph{weight}
each by the likelihood of $o_t$ under its state, $o{=}s\eta,\
\eta\sim\mathrm{LogNormal}(0,\sobs^2)$, so particles consistent with the survey score
high; (iii) \emph{resample} in proportion to weight. The particle cloud \emph{is} the
belief (mean $=$ estimate, spread $=$ uncertainty), and all methods plan on this shared
filter. Particles (not a Kalman filter) are needed because the maps are nonlinear and
the belief is multi-modal (e.g.\ above vs.\ below the Allee threshold).
Under known population identity, $\rbase$ and the $\lambda$ profile are \emph{known
context}: drop $\rbase$ from the latent (do not infer it in the default setting), and
$\reff,\Keff$ are deterministic once $(p,a,\kappa)$ are fixed. Particles thus carry
only $s$ and hidden $(C,\theta,z)$; the log-normal emission is unchanged. The seven
methods only need to read the new cost and the (already known) $(\reff,\Keff)$, and
must \emph{not} read the realized reward (E6 leakage guard).

\paragraph{(E9) Recalibration and declining baselines (do not skip).}
\begin{enumerate}[leftmargin=1.8em,nosep]
\item Re-tune the behavior-policy mixture and per-population calibration profiles so
the healthy-start incident collapse rate returns to $[0.15,0.24]$ \emph{for the
recoverable populations}.
\item \textbf{Real baselines decline.} Four populations have $\lambda_{a_0}<1$
(Egyptian vulture, bottlenose dolphin, Amur tiger, spotted turtle): they decline
under Do-Nothing, so persistence \emph{requires} action --- the decision problem is
``which recovery action, and when.''
\item \textbf{Two are demographic sinks} ($\rmax<0$: Egyptian vulture, bottlenose
dolphin): even $a_4$ leaves $\lambda<1$, so no vital-rate action prevents decline;
only repeated translocation ($a_{10}$) adds individuals. Treat these as hard cases:
report per-population outcomes and split metrics by recoverability rather than
pooling. Optionally exclude sinks from the headline ``beats-both'' statistic and
report them separately.
\end{enumerate}

\paragraph{(E10) Acceptance tests.}
\begin{enumerate}[leftmargin=1.8em,nosep]
\item Env step reproduces \code{action\_effects\_long.csv} for every (population,
action): $\reff,\Delta K,\Delta N,\mathrm{cost}$ match (unit test).
\item For each population, abundance term $\approx0.5$ when $s{=}\Kbase$; $\reff$
within caps; $\Keff\le\Kmax$.
\item Collapse rate in $[0.15,0.24]$ for recoverable populations after recalibration;
document sinks separately.
\item \textbf{Two reward settings:} \code{reward\_mode}=\code{yield} forces
$P_{\mathrm{yield}}{=}0$, \code{safe} uses $P_{\mathrm{safe}}{>}0$; both log identical
benefit ($s_{t+1}$-based, never $o_t$) and cost; the evaluation battery (collapse rate,
min abundance, cost, persistence) is computed independently of \code{reward\_mode} so the
settings stay comparable.
\item \textbf{Reward-leakage guard:} the policy/belief input carries no
realized-reward channel --- two rollouts with identical $(o,a)$ histories but different
logged $R$ must produce identical actions (E6 leakage guard).
\item \textbf{Training convergence (E12):} every agent's logged training loss decreases
and plateaus (no divergence), each W\&B run has a local mirror, and the held-out
validation value is logged and tracks the training curve without diverging (overfitting
check).
\end{enumerate}

\paragraph{(E11) Offline dataset: size convention and behavior-policy diversity.}
\begin{enumerate}[leftmargin=1.8em,nosep]
\item \textbf{$75$k is our convention, not baseline provenance.} The
$\sim\!75$k-transition offline buffer is the Tier-2/3 deep offline-RL convention
(D4RL-style), retained for comparability across the general learners --- \emph{not}
inherited from the ecology baselines. MOOR \cite{moor2025} fits its mechanistic model
from a single short effort/catch series ($T{=}50$ steps, three datasets per
environment) and studies complete-vs-incomplete data, not a transition count; PLUS
\cite{plus2018} is Bayesian adaptive planning over a belief across candidate models
with sequential yearly observations and no fixed buffer. Do not cite PLUS/MOOR to
justify $75$k.
\item \textbf{Optional sample-efficiency ablation.} Sweep the buffer size (e.g.\
$5$k/$25$k/$75$k/$150$k) as a \emph{separate} axis: the deep learners are data-hungry
while the mechanistic baselines are near data-free, so this is an informative
contribution, not a replication.
\item \textbf{Behavior-policy action diversity matters.} Ensure the data-collecting
mixture exercises the action set richly. MOOR found insufficiently variable effort
caused identifiability failures \emph{even with data present}; the analogue here is a
behavior policy that under-uses rate/capacity/translocation actions, leaving $(r,K,N)$
effects unidentifiable. Log per-action coverage as a data-quality check.
\end{enumerate}

\paragraph{(E12) Training instrumentation: experiment logging and a validation split.}
These curves are the first check that a method \emph{actually trained} --- the Tier-3
diagnostic could not tell a failed fit from a hard problem.
\begin{enumerate}[leftmargin=1.8em,nosep]
\item \textbf{Log every run to Weights \& Biases \emph{and} mirror locally.} For each
trained agent (per method $\times$ \code{reward\_mode} $\times$ $\sobs$ $\times$ seed),
log the \emph{training-loss curve} vs.\ gradient step/epoch --- it should decrease and
\emph{converge}. What the curve is depends on the method: the shared $15$-member ridge
dynamics ensemble logs its predictive loss (MSE / negative log-likelihood of
$s_{t+1}$); policy-learning methods (MOPO, OGSRL, RefPlan, Delphic, and MOOR's policy
step) additionally log actor/critic losses ($Q$/Bellman loss, policy loss, and any
penalty or uncertainty term); planning methods (BA-MCTS, PLUS) log the model-fit loss
and rollout diagnostics rather than a policy loss. Also log the run config (method,
reward mode, $\sobs$, population set, seed, buffer size, $\alpha$, $P_{\mathrm{safe}}$),
gradient norms, and the reward decomposition (benefit / cost / penalty). \emph{Mirror
everything locally} (e.g.\ \code{WANDB\_MODE=offline} plus a CSV/Parquet of the metric
history and checkpoints under \code{results/method/reward\_mode/sigma/seed/}) so the
training-error line is recoverable without the cloud.
\item \textbf{Hold out a validation set for a validation curve during training.} Split
the offline buffer \emph{by episode/trajectory} (not by transition, to avoid leakage
across correlated steps), stratified by population and $\sobs$; reserve
$\approx10$--$20\%$ as validation, never used for gradient updates. During training log
a \emph{validation value} on this held-out data: (i) the held-out dynamics predictive
loss ($1$-step and $k$-step) for the ensemble, and (ii) an off-policy value estimate
(e.g.\ FQE or a model-based rollout return) as the policy's validation return. Use it
to watch for overfitting, for early stopping, and for checkpoint/model selection.
\item \textbf{Keep it distinct from the online battery.} This offline validation (on
logged tuples) is \emph{not} the reward-agnostic evaluation battery of E6$'$/E10, which
rolls policies out in the true environment. Compute the validation value under the
\emph{same} \code{reward\_mode} the agent trains on, and --- per the E6 leakage guard
--- reward stays a learning/scoring target only, never a policy or belief-filter input.
\end{enumerate}

\paragraph{Suggested order.} E0 $\to$ E1 $\to$ E2--E5 (transition + cost) $\to$
E6--E6$'$--E7 (state reward + rescale + two settings + conventions) $\to$
E9 (recalibrate) $\to$ E8 (filter) $\to$ E11 (offline data) $\to$ E12 (logging) $\to$
E10 (tests).

\begin{thebibliography}{9}
\bibitem{ricker1954} Ricker, W.E.\ (1954). Stock and recruitment. \emph{J.\ Fish.\ Res.\ Board Can.} 11(5):559--623.
\bibitem{white2022} White, T.B., Petrovan, S.O., Christie, A.P., Martin, P.A.\ \& Sutherland, W.J.\ (2022). What is the Price of Conservation? \emph{BioScience} 72(5):461--471. doi:10.1093/biosci/biac007.
\bibitem{iaconaportal} Iacona, G.\ et al.\ \emph{Conservation Costing Portal} (\code{CostPortal\_7\_26\_24.xlsx}). \url{https://jcleme19.wixsite.com/costingportal/cost-studies}.
\bibitem{plus2018} Memarzadeh, M.\ \& Boettiger, C.\ (2018). Adaptive management of ecological systems under partial observability. \emph{Biological Conservation} 224:9--15.
\bibitem{moor2025} Ju, H.\ et al.\ (2025). Mechanistic model-based offline RL for fishery management. \emph{Expert Systems}.
\bibitem{mopo2020} Yu, T.\ et al.\ (2020). MOPO: Model-based Offline Policy Optimization. \emph{NeurIPS}.
\bibitem{refplan2025} Reflect-then-Plan: Offline Model-Based Planning through a Doubly Bayesian Lens. \emph{ICML} (2025).
\bibitem{bamcts2026} Bayes-Adaptive Monte Carlo Tree Search for Offline Model-based RL. \emph{ICLR} (2026).
\bibitem{ogsrl2025} Yan, R.\ et al.\ (2025). Offline Guarded Safe RL for Medical Treatment Optimization. \emph{NeurIPS}.
\bibitem{delphic2024} Pace, A.\ et al.\ (2024). Delphic Offline RL under Nonidentifiable Hidden Confounding. \emph{ICLR}.
\end{thebibliography}
\end{document}
